\section{展开示例：从床上到书桌前}

\begin{example}[状态画像分析，而不是伪计算]
考虑从躺在床上刷手机，转变为坐到书桌前阅读。
\begin{enumerate}[nosep]
  \item $\Reward(s_0)$ 很高，因为刷手机能持续提供快速的新奇刺激。
  \item $\Switch(s_0)$ 很高，因为读者必须重新找回任务上下文、定位材料，并定义第一个动作。
  \item $\Env(s_0)$ 很高，因为姿势和地点都会一起稳定“继续休息”的行为。
  \item 如果学习任务没有预先指定，$U(s_1)$ 就会很高。
  \item 当书桌被提前整理好、文本事先翻开、并且第一个微任务已经写下时，$P$ 就会上升。
\end{enumerate}
这里的模型并不会计算一个真实的数值障碍。它给出的是一种结构化解释：为什么这个转变会显得昂贵，以及为什么准备能够在行动发生之前就先降低这种昂贵感。
\end{example}

\begin{discussion}
这个例子刻意保持克制。它并没有展示预测成功、参数恢复或实证拟合。它的目的更窄：展示这套记号如何拆开几个原本会被压扁在一起的抱怨，例如“我就是不想做”“我忘了从哪里开始”“这个房间一直把我拉回休息状态”。
\end{discussion}